% Direct manuscript audit; original university scans visually inspected.
\subsection{The cocycle, its primitive, and the manuscript's splittings}

\begin{statement}{D.1}{The manuscript dictionary, with inversion and normalization}{Proved}
Let $F=\widehat{F}(x,y)$ be the geometric profinite tripod group,
with $x=l_0$, $y=l_1$, and $l_\infty=(xy)^{-1}$.
Let $\rho$ be the continuous order-three automorphism
$\rho(x)=y$, $\rho(y)=(xy)^{-1}$. Suppose a continuous
peripheral automorphism $u$ and $g\in F$ satisfy
\[
 u(x)=x^p,\qquad
 u\rho u^{-1}=\operatorname{int}(g)\rho,
 \qquad p\in\widehat\Z^{\times}.
 \tag{D1.1}
\]
Write $g^{\rm ab}=(a_0,a_1)\in\widehat\Z^2$.
Then the raw conjugator $f_0=g^{-1}$ satisfies
$u(y)=f_0^{-1}y^p f_0$. The commutator-normalized
GT-style conjugator is instead
\[
 f=y^{a_1}g^{-1}x^{a_0}\in F',\qquad
 u'=\operatorname{int}(x^{-a_0})u,\qquad
 u'(x)=x^p,\quad u'(y)=f^{-1}y^p f.
 \tag{D1.2}
\]
Here ``GT-style'' describes the normalized pair;
it does not assert the pentagon for an arbitrary
peripheral automorphism.
The change of representative
$u\mapsto\operatorname{int}(x^a)u$ changes
$g\mapsto x^a g y^{-a}$. The element $f$ in
(D1.2) is invariant under this gauge, and
\[
 a_0+a_1=\frac{1-p}{2}.
 \tag{D1.3}
\]

In cote 145, scans 7--12, the manuscript's
$\gamma$ is the orbit of $g$ under precisely
this gauge, and $\chi(u)=p$. Its pair
$(\gamma,\chi)$ therefore corresponds to
$(f,p)$ by (D1.2); it is not literally
$([f],p)$ in the original orbit convention.
Before commutator normalization, inversion
identifies $[g]$ with $[f_0]$ under the opposite
gauge $f_0\mapsto y^a f_0x^{-a}$.
On the manuscript's domain, where the outer
class also commutes with the transposition
$\sigma_\infty$, its additional function $\beta$
is determined by the uncancelled relation
\[
 \sigma_\infty(g^{-1})=x^\beta g y^\beta,
 \qquad \sigma_\infty(x)=y,\quad
 \sigma_\infty(y)=x,
\]
and equals
\[
 \beta=-a_0-a_1=\frac{p-1}{2}.
 \tag{D1.4}
\]
This is a cyclotomic coordinate, not an
identified coordinate of a frame.

Let $\mathcal E$ be the profinite extension
of $S_3\times C_2$ by $F$ associated with
its action on the complex tripod, where
$C_2$ acts by complex conjugation. For a
finite isotropy subgroup $H$ with a
chosen splitting $j:H\to\mathcal E$, a frame
means a geometric conjugator $A$ carrying
$j$ to the transported splitting, with the
specified action on $H$. Its ambiguity is
right multiplication by
$C_{\mathcal E}(j(H))\cap F$. In particular,
if $u(\rho)=A\rho A^{-1}$ in the relevant
extension, then
\[
 g=A\rho(A)^{-1}.
 \tag{D1.5}
\]
Thus the cocycle is the multiplicative
coboundary of a frame, not the frame itself.

The canonical splittings used by LM7 D.25
and LM7\_PUSH E.65 have the stated manuscript
origin: scan 22, equations (11)--(12), gives
the sign-graph $S_3$ isotropy at the
equianharmonic point $P$ and its canonical
lift; scans 27--29 give the isotropy splittings
at $2,-1,1/2$, named $0',1',\infty'$.
Their orientation-preserving subgroups are
respectively the elliptic groups of orders
three and two. These identifications have
now been checked directly against the
original university scans
\citep[scans 7--12, 22, 27--29]{RootGrothendieckCote145}.
They no longer depend on the supplied
AI-made note's reading.
\end{statement}
\begin{proof}
Applying $u\rho=\operatorname{int}(g)\rho u$
to $x$ gives $u(y)=g y^p g^{-1}$.
A conjugator for this one equation is
ambiguous by a right power of $y$ on the
$g$ side. After the specified inner change,
its value is $g'=x^{-a_0}g y^{a_0}$.
Multiplying $g'$ on the right by
$y^{-(a_0+a_1)}$ changes no conjugation
of $y$ and gives $x^{-a_0}g y^{-a_1}=f^{-1}$.
Its abelianization is zero, proving
(D1.2). Under the manuscript gauge,
$a_0$ changes to $a_0+a$ and $a_1$ to
$a_1-a$; substitution proves literal
invariance of the normalized $f$.

The value in (D1.3) follows already in
the continuous quotient $F/\gamma_3F$.
Put $c=[x,y]$ and use its unique class-two
normal forms $x^r y^s c^t$. Multiplication
adds the central correction $-r's$.
Since
\[
 \rho(y)=x^{-1}y^{-1}c^{-1},\quad
 u(y)=y^p c^{a_0p},\quad u(c)=c^{p^2},
\]
the two sides of the intertwining equation,
evaluated on $y$, have normal forms
\[
 \begin{aligned}
 u\rho(y)&=x^{-p}y^{-p}c^{-a_0p-p^2},\\
 \operatorname{int}(g)\rho u(y)
   &=x^{-p}y^{-p}c^{a_1p-p(p+1)/2}.
 \end{aligned}
\]
Their equality gives (D1.3), since $p$
is a unit. This computation is integral
also at two: the binomial expressions are
continuous integral functions on
$\widehat\Z$, and $p$ is odd. Equivalently,
one may multiply the equation by two
and use the absence of additive torsion
in $\widehat\Z$.
Abelianizing the uncancelled scan-12
relation gives $\beta=-a_0-a_1$, hence
(D1.4). The preceding scans independently
write $g y^q\sigma_\infty(g)=x^r$;
its abelianization gives $q=-r$ and
$r=a_0+a_1$, agreeing with that reading.
For the manuscript's complex-conjugation
example, $p=-1$ and $g=y$. Formula (D1.2)
then gives the usual $f=1$, while
$\beta=-1$. This example makes the
need for normalization concrete.

In $F\rtimes\langle\rho\rangle$,
$A\rho A^{-1}=A\rho(A)^{-1}\rho$,
which proves (D1.5). Two conjugators
transporting the same splitting differ
on the right by its centralizer; this
proves the typed ambiguity. The same
argument applies separately to each
finite isotropy subgroup. Concretely, before profinite completion,
$\mathcal E$ consists of pairs $(h,\delta)$
with $h\in S_3\times C_2$ and a homotopy
class of paths from the base point to
its image under $h$, with path composition
as product. A fixed point gives constant
paths and therefore a splitting over its
stabilizer. Moving the base point along
a chosen path transports this splitting;
a different path changes it by geometric
conjugacy. Thus canonical here is literal
at that fixed base point and canonical
up to the stated conjugacy after transport.
At $P$, scan 22
writes the sign graph
$\sigma\mapsto(\sigma,\operatorname{sgn}\sigma)$
inside $S_3\times C_2$ and its canonical
lift. The harmonic discussion identifies
the transposition fixing each of
$2,-1,1/2$ together with complex
conjugation and their lifts. These
are exactly the splittings whose
transports require the respective
order-three and order-two frames.
A finite image of the cocycle alone
need not determine a finite image
of its primitive; no such implication
is used here.

For reproducible provenance, scan numbers
here mean one-based pages of the official
38-page PDF, whose first page is the
archive cover. The printed manuscript
and archive page numbers on the relevant
leaves are one less. The cover identifies
cote 145 as the spring-1981 old redaction.
The retrieval digest and independent
reading record are preserved. The
crossed-out lower argument on scan 12
is not used. In particular, we correct
the supplied note's literal pair
identification by the explicit inversion
and normalization above, rather than
silently adopting its convention.
\end{proof}

\subsection{The questions of Sections 51.15 and 51.16}

The quotations below were first located in the supplied AI-made
transcription and have now been checked against the original official
cote 140-4 images \citep[scans 167--178]{GrothendieckLongueMarche}.
Two independent readings agree on the questions and formulas used here.
Scan numbers mean one-based positions in the 280-page PDF: its first
page is the archive cover, and its printed archive footers are one less.
The cover dates these 279 manuscript pages only as \emph{s.d.}, or
undated; we do not infer a year from cote 145's separately dated cover.
The visible manuscript folio markers place scans 171, 173, 175
and 177 at folios 733, 734, 735 and 736 respectively; hence the intervening
scans are their versos. In particular, scan 178 is folio 736v.
The prompt's shorthand ``733v--735, scans 172--178'' does not
cover the last two marked folios.
The original (357) labels the Tate module of $\mathbb G_m$
by the older twist convention $(-1)$, retained in (358)--(365).
Our modern convention calls the same module $\widehat\Z(1)$.
D.3--D.4 translate that convention; the underlying homomorphisms on
the cyclotomic kernel are unchanged.

\begin{statement}{D.2}{The answer beside each original question}{Proved as a provenance comparison}
The following passages state the original questions answered by the
specified results of Parts A and B. Each answer has exactly the field,
quotient and local hypotheses of those results; a quotient calculation
does not answer the full extension question.
\end{statement}
\begin{proof}
In Section 51.16, folio 733, scan 171, equation (366), the curve is the
\emph{plus} curve $x_0^n+x_1^n+x_\infty^n=0$ with all coordinates nonzero.
The original says that it ``n'a pas de pt rat $/\mathbb Q$''.
For odd $n$ the substitution $z=-x_\infty$ identifies it with the
classical curve. For even $n$ the plus curve has no real point at all,
whereas the classical curve has four rational cusps. Statements A.1--A.16,
A.20--A.21 and A.31 give the section answer, keeping this distinction:
the odd full Selmer section set is the union of the three rational cusp
packets; the classical even answer is the union of four packets; the plus
even answer is empty. The assertions for all sections use the named local
Theorem 1.1. The local calculations A.3--A.6 and A.28--A.30 identify the
arithmetic inputs needed for this answer rather than adding a new reading
to the manuscript.

In Section 51.16, folio 733v, scan 172, equation (371), the text distinguishes
the splittings that are ``distincts des `scindages \`a l'infini' ''.
On folio 734, scan 173, it asks whether one can prove
``l'inexistence de scindages de (371) autres que ceux `provenant des pts
\`a l'infini' ''. The same page proposes ``des `approximations' de plus
en plus fines de $\Pi_1(U_n)$ et de la structure d'extension (371)''.
The repaired full Selmer theorem in A.1--A.16 answers the first question
for the Fermat curves in question, and A.20--A.21 give the exact additional
local input for all full sections. The twist results A.23--A.27 and A.32--A.43
extend that answer only in their stated rank-zero classes. They retain
the infinitely many classes in each cusp packet, a datum absent from a
mere list of cusp locations.

The quotient question is explicit in Section 51.16, folio 735, scan 175,
equation (373), and folios 735v--736, scans 176--177, equations (377)--(379).
On scan 177 the text asks to ``d\'eterminer toutes ces classes de
$\widehat M[n]$-conjugaison possibles de tels scindages'' and then asks
which can correspond to a rational point. Statements B.1--B.22 compare
the full, relative prosolvable, maximal prosolvable, derived, pro-$3$
nilpotent and metabelian problems with this question. The subgroup and
coefficient distinctions in those statements are essential: the cubic
abelian Selmer result does not identify all abelian sections, and the
infinitely ramified class-three construction is not a full profinite
section. The general intrinsic action, arbitrary-input metabelian
classification and different-prime colouring remain unanswered parts
of this original request for finer approximations. B.15--B.17 give the requested explicit
phantom coordinates and the fibre comparison. B.19--B.22 separate the
Hodge regulator from the arithmetic action and correct the scope of its
proposed integral test. B.35--B.40 subsequently determine the arithmetic
Ross component and its actual integral refinement, with the intrinsic
finite detector equal to $19$ modulo $27$.
B.41 proves that this extension alone does not determine arbitrary
class-three lifting: actual class-two cusp origins supply additional
arithmetic data. B.42 proves a common cusp colour at the specified inert
primes while retaining the unrestricted two-prime question as open.
A.69 determines the actual joint level-three lifting exponents for
every finite alteration of every integer on the least cubic; it does
not infer a classification of arbitrary infinite alterations.

Finally Section 51.16, folio 736v, scan 178, says
``l'approximation sur laquelle nous travaillons est beaucoup trop
grossi\`ere'' and concludes that ``il est n\'ecessaire de passer \`a des
approximations plus fines''. B.2's relative prosolvable quotient gives one
precise refinement that carries the Fermat theorem. B.7 and B.15--B.17
show why the existence of a section in a nilpotent
or metabelian approximation alone does not give a section of (371).
No extrapolation along a noncofinal tower is used in this comparison.

For Section 51.15, folio 731v, scan 168, questions (358)(a)--(b) ask whether
the Kummer maps attached to distinct primes are surjective or have open
image, separately and jointly. Folio 732, scan 169, decomposes them into
coefficient primes; folio 732v, scan 170, asks the stronger local-inertia
question (365). A.29 and A.32 use exactly the completed Kummer parameters
whose components this passage distinguishes. D.3 below answers the joint
global image question. It does not claim to compute every local-inertia
image in (365). This is also the historical source of the coefficient and
completion distinctions retained throughout Part B.
Statement D.4 below subsequently determines that local-inertia image
itself, including its exceptional factor at the coefficient prime two.

There is also a mathematical inconsistency in the original scan 178,
confirmed directly rather than attributed to the transcription: the displayed
affine equation there has right side $1$, whereas equation (366) defines
the plus curve by the homogeneous sum equal to $0$. We use (366) and
the actual coordinate-ratio embedding in the torus, as proved in A.17;
we do not use the inconsistent right side to identify a cover or its
arithmetic fundamental group.
\end{proof}

\begin{statement}{D.3}{The joint global Kummer image in Section 51.15}{Proved}
Let $p_1,\ldots,p_m$ be distinct positive rational primes and put
$G_{\Q}^{+}=\Gal(\overline\Q/\Q^{\mathrm{ab}})=\ker\chi$.
Fix compatible roots and compatible roots of unity. The joint Kummer map
\[
G_{\Q}^{+}\longrightarrow\wh\Z(1)^m
\]
has image $2\wh\Z(1)^m$. Thus it has index $2^m$, is surjective at every
odd coefficient prime, and has image $2\Z_2(1)^m$ at $2$. In particular,
the surjectivity alternative in Section 51.15, scan 168, is false, while
its open-image alternative holds with this exact index. This concerns
the full cyclotomic kernel, not the smaller local inertia group of (365).
\end{statement}
\begin{proof}
All square roots of rational numbers lie in $\Q^{\mathrm{ab}}$, so every
Kummer coordinate is even. To prove that these are the only restrictions,
work at a finite level $n$ over the field $K=\Q^{\mathrm{ab}}$, which
contains $\mu_n$. Kummer duality says that the annihilator of the image
in $(\Z/n)^m$ consists of the exponent vectors $(e_i)$ for which
$a=\prod_i p_i^{e_i}$ has an $n$th root in $K$.

Choose the positive real root $\alpha=a^{1/n}$. If $\alpha\in K$, complex
conjugation commutes on $K$ with every $\sigma\in G_\Q$. Writing
$\sigma(\alpha)=\zeta\alpha$ gives both
$c\sigma(\alpha)=\zeta^{-1}\alpha$ and
$\sigma c(\alpha)=\zeta\alpha$. Therefore $\zeta^2=1$ for every
$\sigma$, so $\alpha^2\in\Q$. Unique factorization in $\Q^\times$ now
gives $n\mid2e_i$ for every $i$. Conversely, those divisibilities make
$\alpha^2$ rational, so $\alpha$ is quadratic over $\Q$ and belongs to
$K$. The annihilator is therefore exactly
\[
\{(e_i)\in(\Z/n)^m:2e_i=0\ \text{for every }i\}.
\]
Its orthogonal subgroup is $2(\Z/n)^m$. The image of the compact group
$G_\Q^+$ is closed in $\wh\Z^m$, and its reductions have now been
determined for every $n$. It must equal $2\wh\Z^m$. This proof also
shows independence of the selected primes; checking their images one at
a time would not suffice for the joint assertion.
\end{proof}

\IfFileExists{tables/manuscript_concordance.tex}{\subsubsection*{Individual concordance for Parts A and B}
The short question keys in the statements and the table refer to the
following passages, checked against the original official manuscript images
after being located in the supplied transcription. \textbf{F}
asks for ``l'inexistence de scindages de (371) autres que ceux
`provenant des pts \`a l'infini' '' (Section 51.16, folio 734, scan 173).
\textbf{A} asks to ``d\'eterminer toutes ces classes de
$\widehat M[n]$-conjugaison possibles de tels scindages'' (Section 51.16,
folio 736, scan 177). \textbf{T} refers to that same question for the
explicit extension (373), constructed on folios 735--736, scans 175--177.
\textbf{K} asks ``sont-ils surjectifs, ou du moins l'image est-elle un
sous-groupe ouvert ?'' and asks this jointly for distinct primes
(Section 51.15, folio 731v, scan 168, (358)(a)--(b)). \textbf{L} asks
``si cet homomorphisme est surjectif, ou si du moins son image est
ouverte ?'' for the restriction to $I_\ell^+$ (Section 51.15,
folio 732v, scan 170, (365)).

The table states the answer or the exact technical contribution of
\emph{each} statement. Intermediate lemmas are not represented as new
questions literally written in the manuscript. The qualifications in
each numbered statement, including the distinction between full and
quotient sections, are retained.
\begingroup\small\setlength{\tabcolsep}{4pt}
\begin{longtable}{@{}p{.09\textwidth}p{.10\textwidth}p{.75\textwidth}@{}}
\toprule ID & Question & Answer and status \\ \midrule
\endfirsthead
\toprule ID & Question & Answer and status \\ \midrule
\endhead
A.1 & F & \textbf{Known}. Full all-place Selmer sections stay Selmer in finite étale neighbourhoods; compactifications yield adelic points surviving every finite torsor. \\
A.2 & F & \textbf{Proved}. Full Selmer plus finite support on a density-one set of finite places implies a point or rational cusp section; the target finite-descent equality consequence holds. \\
A.3 & F & \textbf{Proved}. For q=a u\textasciicircum{}e h(u), the cusp-packet map is b maps to kappa(a)+eb; each Fermat--Legendre cusp has q-class -8kappa(2)+2ellkappa(u). \\
A.4 & F & \textbf{Refuted}. The asserted cusp divisibility by ell in LM6 E.8 and Steklov's general ramification-divisibility proposition is false. \\
A.5 & F & \textbf{Proved}. Corrected local quadratic twists are unramified outside 2ell, finite flat/weight2 at ell, and tame of conductor exponent1 at2, for every point and cusp case. \\
A.6 & F & \textbf{Known}. The mod-ell Legendre representation of a Selmer section of Uell is reducible and has a quadratic twist with semisimplification 1 plus chiell; the irreducibility test has level2. \\
A.7 & F & \textbf{Proved}. The cyclic-ell Legendre cover maps nonconstantly over Q to Y0(4ell); section-supporting components are geometric and g(X0(4ell))=(ell-3)/2. \\
A.8 & F & \textbf{Proved}. Every full profinite Selmer section of Uell/Q is cuspidal for every prime ell\textgreater{}=5. \\
A.9 & F & \textbf{Refuted}. Y1(7)(Q) is nonempty; y\textasciicircum{}2-xy-4y=x\textasciicircum{}3-4x\textasciicircum{}2 has discriminant-1664 and (0,0) of exact order7. \\
A.10 & F & \textbf{Proved}. Every full profinite Selmer section of U3/Q is cuspidal; its compactification is X0(27) and has exactly the three rational cusps. \\
A.11 & F & \textbf{Proved}. Every full profinite Selmer section of the plus Fermat curve Un/Q is cuspidal for all odd n\textgreater{}=3, with exactly three rational cusp centres. \\
A.12 & F & \textbf{Proved}. The three classical quartic quotients are E-:y\textasciicircum{}2=x\textasciicircum{}3-4x and twice E+:y\textasciicircum{}2=x\textasciicircum{}3+4x; all have geometric CM ring Z[i], and Jac(C4) is Q-isogenous to E-×E+\textasciicircum{}2. \\
A.13 & F & \textbf{Proved}. Exact Mordell--Weil groups are E-(Q)=C2×C2 and E+(Q)=C4, with complete explicit point lists. \\
A.14 & F & \textbf{Proved}. Exact eta-product positivity proves L(E27,1), L(E-,1), L(E+,1)\textgreater{}0, and Rubin gives finite Tate--Shafarevich groups. \\
A.15 & F & \textbf{Proved}. Each of the three elliptic quotients separately proves every full profinite Selmer section of classical C4/Q is one of its four rational cusp packets. \\
A.16 & F & \textbf{Proved}. For every n\textgreater{}=3, full profinite Selmer sections of classical Cn/Q are precisely images of its rational cusp packets, three centres for odd n and four for even n. \\
A.17 & T, A & \textbf{Proved}. The torus quotient has the full two-coordinate completed Kummer section set; its S3-equivariant splitting is unique up to unique geometric conjugation. \\
A.20 & F & \textbf{Proved}. The stipulated proper local Theorem 1.1 implies the affine point-or-cusp alternative by cofinal neighbourhoods; only this named additional input promotes the global Selmer theorem to all sections. \\
A.21 & F & \textbf{Known}. Real sections are classified by connected components. Two cusp sides can give the same component; the plus curve of even exponent has no real section. \\
A.22 & F & \textbf{Proved}. A density-one support assertion cannot replace all-place descent. The displayed density-intersection inference fails; its repair uses strong positive upper density. \\
A.23 & T, A & \textbf{Proved}. The completed Kummer fibres of each cusp packet are computed exactly. Whenever a packet is present the set of section classes is infinite. \\
A.24 & F & \textbf{Proved}. A nonconstant map to a rank-zero target whose Tate--Shafarevich group has no nonzero divisible subgroup gives only rational-point or rational-cusp sections. \\
A.25 & F & \textbf{Proved}. The three elliptic quotients of every diagonal quartic are written over Q, with their exact coefficients and CM by Z[i]. \\
A.26 & F & \textbf{Proved}. Explicit quotient maps give rank-zero criteria for diagonal sextics; the complete ten-factor decomposition is supplied in A.34. \\
A.27 & F & \textbf{Proved}. The diagonal cubic Jacobian is y\textasciicircum{}2=x\textasciicircum{}3-432(abc)\textasciicircum{}2; its isomorphism class depends on the cube class of abc. \\
A.28 & F, L & \textbf{Proved}. At good primes, a ramified section in the full prime-to-residue-characteristic quotient lies at one cusp; this uses a single simultaneous quotient. \\
A.29 & K, F & \textbf{Proved}. Cross-prime agreement concerns one completed Kummer parameter, not one ordinary integer. An explicit parameter with different primary components refutes the latter reading. \\
A.30 & F, L & \textbf{Open}. The exact independent weakest local hypothesis remains open. The residual modularity and all-place descent requirements used by the proof are separately stated. \\
A.31 & F & \textbf{Proved}. With the stipulated local theorem, the full global section is cuspidal, which supplies common cusp and parameter data at every prime. \\
A.32 & F, K & \textbf{Proved}. The rank-zero fibres exclude noncuspidal ghosts; cusp Kummer parameters can have infinite ramification, so residual finite ramification is insufficient. \\
A.33 & F & \textbf{Proved}. Elementary torsion classifications bound the two CM families by orders six and four and list every nonzero-coordinate torsion case used by the point census. \\
A.34 & F & \textbf{Proved}. All ten elliptic factors of a diagonal sextic are supplied, with maps and an independent holomorphic-differential basis. \\
A.35 & F & \textbf{Proved}. Universal lists of three tripod values for cubic and quartic criteria, and fifteen for sextic criteria, give every ordinary point in the settled classes. \\
A.36 & F & \textbf{Proved}. For every class meeting a CM rank-zero criterion, the section set is the explicit ordinary-point list together with all rational cusp packets. \\
A.37 & F & \textbf{Proved}. Every analytic-rank label in the cubic, quartic and sextic coefficient tables has an exact root-number or non-torsion-point vanishing proof and a rigorous Arb nonvanishing interval with an explicit infinite-tail and Taylor-remainder bound. \\
A.38 & F & \textbf{Proved}. The 81 positive cube-free coefficient classes supported at 2,3,5,7 have analytic-rank counts 34,41,6 in ranks 0,1,2 respectively, with no unresolved classes. \\
A.39 & F & \textbf{Proved}. All 6561 cubic pairs are tabulated; the rank-zero criterion settles 2754, with 2654 empty section sets, 97 cusp-only cases and 3 ordinary-plus-cusp cases. \\
A.40 & F & \textbf{Proved}. The 512 quartic coefficient curves y\textasciicircum{}2=x\textasciicircum{}3+Ax supported at 2,3,5,7 have exact analytic-rank counts 220,256,36 in ranks 0,1,2 respectively. \\
A.41 & F & \textbf{Proved}. The 2592 sextic coefficient curves y\textasciicircum{}2=x\textasciicircum{}3+B supported at 2,3,5,7 have exact analytic-rank counts 1136,1296,160 in ranks 0,1,2 respectively. \\
A.42 & F & \textbf{Proved}. The complete quartic 262144-row and sextic 6718464-row CSVs record every quotient criterion and every Selmer section in the settled classes, using exact class, point and cusp-packet dictionaries. \\
A.43 & F & \textbf{Proved}. The smallest conductor among analytic-rank-one cubic twists y\textasciicircum{}2=x\textasciicircum{}3-432d\textasciicircum{}2 is243, attained by d=9, with minimal model y\textasciicircum{}2+y=x\textasciicircum{}3-1 and a certified nonzero first derivative. \\
A.44 & F & \textbf{Proved}. Every profinite multiple of a rational non-torsion elliptic point lifts through every finite-kernel refinement, or finite collection of such refinements. This is not a Selmer assertion. \\
A.45 & F & \textbf{Open}. The full intrinsic metabelian and congruence lifting sets remain open. A59--A69 settle the stated abelian exponent classes and every finite alteration of an integer at the actual joint level-three quotient. \\
A.46 & F, L & \textbf{Proved}. For every full section of U\_ell, the residual Legendre representation becomes unramified after a quadratic twist at every q outside \{2,ell\} union the prime divisors of ell\textasciicircum{}2-1. Local point-or-cusp input only at that finite set suffices for the residual reducibility theorem. \\
A.47 & F, L & \textbf{Proved}. The natural global quotient retaining the entire prime-to-three geometric group of U7 and its residual Legendre7 representation admits a section with the prime-to-three data of a global cusp and wild residual inertia of order3 at3, unremovable by a quadratic twist. \\
A.48 & F, L & \textbf{Proved}. Intrinsic ramification and ramification of a cocycle in a chosen splitting agree at good reduction but cannot be interchanged at bad primes; for q|n odd Fermat boundary inertia already acts nontrivially in geometric abelianization for every section. \\
A.49 & F, L & \textbf{Proved}. The proper elliptic curve X0(27) has a Z5-family of full global sections whose nonzero members are point-induced at every completion except Q5 and are nongeometric at Q5. \\
A.50 & F & \textbf{Known}. The Jacobian J of C:y\textasciicircum{}2+y=x\textasciicircum{}5 has conductor3125 and CM by Q(zeta5), becomes a Hilbert modular abelian surface of level(sqrt5)\textasciicircum{}3 over F=Q(sqrt5), and satisfies L(J/F,s)=L(J/Q,s)\textasciicircum{}2. Nonvanishing at1 implies finite J(F) and finiteSha(J/F) by the precisely stated Hilbert modular rank-zero theorem. \\
A.51 & F & \textbf{Proved}. The root number of Jac(y\textasciicircum{}2+y=x\textasciicircum{}5)/Q is +1 and L(J/Q,1) lies in the certified ball 0.71831362842056477051660347805 ± 2.72e-30. \\
A.52 & F & \textbf{Proved}. Every Selmer section of the geometrically maximal prosolvable arithmetic fundamental group of the plus U5 is cuspidal. \\
A.53 & F & \textbf{Proved}. The entire plus Fermat quintic Jacobian is Q-isogenous to three copies of Jac(y\textasciicircum{}2+y=x\textasciicircum{}5), and therefore has finite rational Mordell--Weil group and finite Tate--Shafarevich group. \\
A.54 & F & \textbf{Proved}. Both Jac(y\textasciicircum{}2+y=x\textasciicircum{}5)(Q5) and Jac(X5)(Q5) have no prime-to-five torsion and are pro-five groups; their full local five-adic Tate Kummer maps are injective. The explicit isogeny and its transpose have kernels killed by five. \\
A.55 & F & \textbf{Proved}. The full local point section at [2:y:1] on U7/Q7, with y\textasciicircum{}7=-129, is intrinsically ramified in the geometric pro-three quotient and is not cuspidal even in geometric abelianization. \\
A.56 & F & \textbf{Proved}. For every prime q dividing ell(ell\textasciicircum{}2-1), the natural arithmetic quotient of pi1(U\_ell) retaining its entire prime-to-q geometric quotient and residual Legendre representation is the full fibre product; any global determinant-cyclotomic residual representation can be paired with a fixed global prime-to-q cusp section. \\
A.57 & F & \textbf{Proved}. At every exceptional prime, actual global elliptic residual representations give global joint-quotient sections with fixed complete prime-to-q cusp data but unavoidable bad local behavior: nonzero inertia for odd q distinct from ell, non-finite-flatness at ell, or conductor exponent at least five at two, even after every quadratic twist. \\
A.58 & F & \textbf{Proved}. A full U\_ell section is cuspidal if its residual Legendre representation admits a quadratic twist unramified outside2ell, finite flat atell and of conductor exponent at2 at most4, and one invariant-line section has a cofinal family of neighbourhood compactifications with nonempty finite-abelian-descent adelic loci. \\
A.59 & F & \textbf{Proved}. Every proper section over Q lifts to the full cuspidally central quotient of an elliptic curve punctured at a union of torsion sub-packets. In particular every profinite rank-one exponent on the conductor243 cubic lifts; its fibre is the product of four completed multiplicative groups specified by the four cusp residue fields. \\
A.60 & F & \textbf{Proved}. For every elliptic curve E/k, full torsion divisor D=E[N], nontorsion R in E(k), and nonempty finite coefficient-prime set P, the proper Kummer class (1-e\_P)kappa(R) admits no geometrically abelian Selmer lift to E-D. On the rank-one cubic this excludes all n(1-e\_P), n a nonzero integer, including the exponent zero at2 and one at all odd primes. \\
A.61 & F & \textbf{Proved}. For a nontorsion point on a CM elliptic curve with all CM endomorphisms defined, the p-adic joint torsion/Kummer image is open for every p, and the joint image for finitely many distinct coefficient primes is open. After every finite base extension a finite Frobenius condition forces all these primary components of the point to vanish locally. All finite quotient prime divisors lie among primes dividing the products p(p\textasciicircum{}2-1). \\
A.62 & F & \textbf{Proved}. For an abelian Selmer exponent equal to one integer outside a finite coefficient-prime set, every global unit character that extends as a unit at that integer point agrees with its rational-point value in all coefficient components outside one explicit finite set. \\
A.63 & F & \textbf{Proved}. On every CM elliptic curve punctured at all N-torsion, an abelian Selmer compact exponent which equals one integer at all but finitely many coefficient primes must equal that integer everywhere. The lifting set on all finite alterations of integers is therefore exactly Z, also for the intrinsic metabelian Selmer quotient. \\
A.64 & F & \textbf{Proved}. For any elliptic curve punctured at at least two geometric points, the full geometric abelian Selmer fibre over a rational compact point is the unique point section when the centre is interior, and the complete injectively parametrized completed-Kummer cusp packet when the centre is a rational boundary point. \\
A.65 & F & \textbf{Proved}. For any elliptic E/Q punctured at E[N], N\textgreater{}=2, and nontorsion R in E(Q), the exponents a in Zhat which are units at all but finitely many coefficient primes admit a full geometric abelian Selmer lift exactly when a is a nonzero integer. Their abelian fibres are unique point sections. The same exponent set holds for the intrinsic metabelian quotient. Over any number field, u(1-e\_P)kappa(R) is excluded for every profinite unit u and nonempty finite P. \\
A.66 & F & \textbf{Proved}. For arbitrary E/Q, U=E-E[N], and nontorsion rational R, let a have finitely many zero coefficient components and suppose some prime lambda divides v\_p(a\_p) for all but finitely many nonzero components. Then abelian Selmer lifting of a*kappa(R) is equivalent to a being a nonzero integer, with unique point-section fibre. The same exponent set holds metabelianly. \\
A.67 & F & \textbf{Proved}. For the actual joint quotient of the full compact cubic group with finite projective j-line modular level m, every finite coefficient erasure n(1-e\_P) fails Selmer lifting at m=3 if3 is not inP, and at m=4 for every finite nonemptyP. On the least rank-one cubic E9, every singleton erasure n(1-e\_p), all primesp and all n!=0, already fails at m=3. \\
A.68 & F & \textbf{Proved}. On the least rank-one cubic, the actual finite joint j-line level-three quotient admits Selmer lifting among every finite alteration of a nonzero integer exactly at that integer. Pure coefficient-three alterations of any integer, including0, are likewise exactly classified. Uniformly for elliptic CM cubic twists with rational infinity flex, any finite alteration of a nonzero integer with a changed prime other than3 is excluded, including coefficient2. \\
A.69 & F & \textbf{Proved}. On E:y\textasciicircum{}2+y=x\textasciicircum{}3-1 with R=(1,0), U=E-E[3], the actual joint compact/j-line level-three quotient admits a Selmer lift of akappa(R), for an exponent a differing from an integer n at only finitely many coefficient primes, if and only if a=n. In particular every nonzero finite-support exponent is excluded, completing the baseline-zero case. \\
A.70 & F & \textbf{Proved}. For every cubic E\_d:X\textasciicircum{}3+Y\textasciicircum{}3=dZ\textasciicircum{}3 with rational infinity origin, nontorsion R in E\_d(Q), and the actual joint compact/j-line level-three quotient, any pure coefficient-three alteration of any integer n has a Selmer lift exactly when it equals n. Together with A.68, this classifies every finite alteration of every nonzero integer uniformly across all these twists, including twists ramified at2. \\
A.71 & F & \textbf{Proved}. For every rational-origin cubic E\_d:X\textasciicircum{}3+Y\textasciicircum{}3=dZ\textasciicircum{}3 with a nontorsion rational point R, and every exponent a differing at finitely many coefficient primes from any integer n, the actual joint compact/j-line X3 quotient admits a Selmer lift of akappa(R) exactly when a=n. This includes arbitrary finite support around0, uniformly across all such twists. \\
B.1 & A, F & \textbf{Proved}. Finite support for the maximal prosolvable quotient \\
B.2 & A, F & \textbf{Proved}. Selmer sections of the Legendre-kernel derived-residual quotient are cuspidal \\
B.3 & A, F & \textbf{Proved}. Cubic and classical quartic maximal prosolvable Selmer sections are cuspidal \\
B.4 & A, F & \textbf{Proved}. Septic maximal prosolvable Selmer sections are cuspidal \\
B.5 & A, F & \textbf{Proved}. Least positive derived length for cubic is k=1 \\
B.6 & A, F & \textbf{Proved}. Every nonzero E31/E33/E34 family member is non-Selmer at its coefficient prime \\
B.7 & A & \textbf{Refuted}. Infinitely ramified noncuspidal class-two sections can lift to class three \\
B.8 & A & \textbf{Known}. Quartic class-two outer action has rank-two CM carrier and Ceresa factor two \\
B.9 & A & \textbf{Proved}. cd3 of Q(mu3-infinity) is at most one; actual full pro3 lifts exist and lattice H2cts vanishes \\
B.10 & A & \textbf{Proved}. Continuous H2 is literal cyclotomic cokernel and phantom is closure/image \\
B.11 & A & \textbf{Proved}. R12 restricted-family algebra and non-Mittag--Leffler defect \\
B.12 & A, F & \textbf{Proved}. Exact Selmer classification of all R9 constructed families and least positive derived lengths \\
B.13 & A & \textbf{Open}. The remaining arbitrary-input class-three and intrinsic metabelian/marked maps \\
B.14 & A & \textbf{Open}. Full profinite metabelian sets, cross-prime cusp colour equality and full pro3 infinite-ramification equivalence \\
B.15 & A & \textbf{Proved}. Explicit 44-coordinate phantom in the cubic class-two lifting problem \\
B.16 & A & \textbf{Proved}. LM7\_PUSH A4 phantom survives the class-two comparison; fibre ranks satisfy 24=23+1 \\
B.17 & A & \textbf{Proved}. LM6 E38 is the same nonclosed-image phantom mechanism in degree three \\
B.18 & A, F & \textbf{Proved}. All cubic E16/E34 families with infinite compact class fail the Selmer condition at three \\
B.19 & A & \textbf{Known}. The Ross symbol has exact real-regulator coordinate -3/2 in Ito's normalized L-value basis \\
B.20 & A & \textbf{Proved}. Exact rational elliptic action comparison with the Ross cup edge, including scalar -1 \\
B.21 & A & \textbf{Proved}. E39 gives the three-dimensional abelian-path kernel; the ambient eta lattice has substantially larger torsion \\
B.22 & A & \textbf{Refuted}. The fixed order-3\textasciicircum{}30 finite Frobenius test does not determine the integral action-extension class \\
B.23 & A & \textbf{Proved}. Finite-base quartic cohomology at 3 and 5, with all division-field and class-group certificates \\
B.24 & A & \textbf{Proved}. Complete unified class-two local-order law for the cubic at 3 and plus quartic at 3 and 5 \\
B.25 & A, F & \textbf{Proved}. Unbounded quartic local Kummer orders and failure of additivity of all three class-two images \\
B.26 & A, F & \textbf{Proved}. Quintic full geometric abelian Selmer sections are cuspidal; least positive derived length is one \\
B.27 & A & \textbf{Proved}. All whole cocycles unramified outside 3 or 5 have exactly three Selmer packets at every corresponding Fermat prime-power level \\
B.28 & A & \textbf{Proved}. Every nonempty intrinsic class-three fibre for every input is explicitly evaluated \\
B.29 & A & \textbf{Proved}. The actual local maximal pro-three Galois groups at the two primes above seven form the free pro-three product G\_\{Q(zeta3),\{3,7\}\}\textasciicircum{}(3) \\
B.30 & A & \textbf{Refuted}. Task B.10(b) is false for the full geometric pro-three group: a Z3-family of sections unramified outside3,7 has nonzero seven-adic inertia at a non-rational K-cusp and is not globally cuspidal \\
B.31 & A & \textbf{Proved}. Entire zero-abelian class-three slice has an exact linear local-order/decay law and unbounded orders \\
B.32 & A & \textbf{Proved}. Inherited tripod metabelian action kills the compact determinant extension defining eta \\
B.33 & A & \textbf{Proved}. Complete section coordinates unramified outside3,7: exact two-variable local relation modulo fixed complex-conjugation group, all local-class fibres, and failure of local geometry at3 for B.30 \\
B.34 & A & \textbf{Refuted}. The rational local colour O at7 also occurs for globally noncuspidal full pro-three sections; explicit path conjugation creates an infinite-order compact class \\
B.35 & A & \textbf{Proved}. Every good-prime determinant extension is nonzero; the local2 problem reduces to an exact affine set of 59049 module classes \\
B.36 & A & \textbf{Proved}. The actual cubic determinant-extension class at2 has exact order27; the coefficient threshold is27 \\
B.37 & A & \textbf{Proved}. The actual outer eta and the specified arithmetic Ross cup-edge are rationally nonzero; the arithmetic space has dimension one at3 \\
B.38 & A & \textbf{Proved}. The homogeneous integral action ambiguity is F3\textasciicircum{}4 under global paths, and C3 after the integral commutator and local quaternion constraints \\
B.39 & A & \textbf{Proved}. An actual Frobenius-stable intrinsic class-two quotient of order3\textasciicircum{}36 and exponent27 detects the remaining integral action phase through an explicit primitive statistic of its outer Frobenius automorphism \\
B.40 & A & \textbf{Proved}. The actual geometric extension selects eta at2 exactly and determines the integral Ross refinement \\
B.41 & A & \textbf{Proved}. The entire Ross extension and geometric outer data do not determine arbitrary-input class-three lifting; the changed cusp origins are explicit \\
B.42 & A & \textbf{Proved}. Every actual full cubic pro-three section has one common rational cusp colour at all ramified inert primes q with v3(q+1) \textless{}= 3+v3(b), where b is its exact free compact cohomology coordinate; the compact twenty-seven bound is sharp; split q=4,7 mod9 satisfy an exact Frattini compact-colour formula \\
B.43 & A & \textbf{Proved}. Two actual cubic cusp origins have compatible explicit Heisenberg projections and evaluated finite Frobenius values \\
\bottomrule
\end{longtable}
\endgroup
}{}

\begin{statement}{D.4}{The exact local-inertia Kummer image}{Proved}
Fix a rational prime $\ell$, a decomposition group at $\ell$, its
inertia subgroup $I_\ell$, and
$I_\ell^+=\ker(\chi_\ell:I_\ell\to\Z_\ell^\times)$.
For a positive rational prime $p$, the image of the continuous Kummer
homomorphism $\kappa_\ell(p)|I_\ell^+$ is
\[
\begin{cases}
 \Z_\ell(1),&\ell\ne2, p=\ell,\\
 \ell^{v_\ell(p^{\ell-1}-1)-1}\Z_\ell(1),&\ell\ne2, p\ne\ell,\\
 2\Z_2(1),&\ell=p=2,\\
 2^{v_2(p^2-1)-2}\Z_2(1),&\ell=2, p\ne2.
\end{cases}
\tag{D4.1}
\]
Thus the open-image question in Section 51.15, folio 732v, scan 170,
equation (365), has an affirmative answer for each $p,\ell$.
Surjectivity is false in general, and always fails at $\ell=2$.

More generally, for $a_1,\ldots,a_m\in\Q_\ell^\times$, write
\[
 a_i=\ell^{v_i}\omega_i(1+\ell)^{d_i}\quad(\ell\ne2),
 \qquad a_i=2^{v_i}\epsilon_i5^{d_i}\quad(\ell=2),
\]
where $d_i\in\Z_\ell$, $\omega_i\in\mu_{\ell-1}$ and
$\epsilon_i\in\{1,-1\}$. Put $v=(v_i)$ and $d=(d_i)$.
After a compatible identification of the Tate twists, the joint image is
\[
 \Z_\ell v+\Z_\ell d\quad(\ell\ne2),
 \qquad 2\Z_2v+2\Z_2d\quad(\ell=2).
\tag{D4.2}
\]
In particular its rank is at most two, and it is not open in
$\Z_\ell(1)^m$ when $m>2$. The joint local assertion is therefore
strictly different from the joint global assertion D.3.
\end{statement}
\begin{proof}
The local Kronecker--Weber theorem identifies
\[
 \Q_\ell^{\mathrm{nr}}(\mu_{\ell^\infty})
   =\Q_\ell^{\mathrm{ab}},
 \qquad I_\ell^+=G_{\Q_\ell^{\mathrm{ab}}};
\]
see \citet[I, Theorem 4.8 and Corollary 4.12]{BMilneCFT2020},
with the multiplicative Lubin--Tate group for $\Q_\ell$.
We first calculate which radicals of an element $a\in\Q_\ell^\times$
can lie in this abelian extension.

If $a^{1/\ell^n}\in\Q_\ell^{\mathrm{ab}}$, its finite Kummer cocycle
$c$ factors through the abelian group $G_{\Q_\ell}^{\mathrm{ab}}$.
The cocycle identity and commutativity give, for every $g,h$,
\[
 (\chi_\ell(g)-1)c(h)=(\chi_\ell(h)-1)c(g)
 \quad\text{in }\Z/\ell^n(1).
\tag{D4.3}
\]
For odd $\ell$, choose $g$ with $\chi_\ell(g)=2$.
Then $c(h)=(\chi_\ell(h)-1)c(g)$ is a coboundary. Kummer theory
implies $a\in(\Q_\ell^\times)^{\ell^n}$; the converse is immediate.
For $\ell=2$, choose $g$ with $\chi_2(g)=-1$. Equation (D4.3)
shows that $2c$ is a coboundary, hence
$a^2\in(\Q_2^\times)^{2^n}$, or equivalently
\[
 a\in\{1,-1\}(\Q_2^\times)^{2^{n-1}}.
\tag{D4.4}
\]
Conversely, if $a=\epsilon b^{2^{n-1}}$, a $2^n$th root is the
product of a square root of $b$ and a $2^n$th root of $\epsilon$.
Both lie in $\Q_2^{\mathrm{ab}}$, so (D4.4) is exact.

At level $\ell^n$, the annihilator of the joint Kummer image consists
of $e\in(\Z/\ell^n)^m$ such that
$\prod_i a_i^{e_i}$ has an $\ell^n$th root in
$\Q_\ell^{\mathrm{ab}}$. For odd $\ell$, the preceding calculation
says precisely $v\cdot e=d\cdot e=0$ modulo $\ell^n$.
For $\ell=2$, it says $2v\cdot e=2d\cdot e=0$ modulo $2^n$;
the signs contribute no condition. Finite-module duality identifies
the orthogonal subgroup with the span of $v,d$, or of $2v,2d$,
respectively. Passing to inverse limits gives (D4.2), since the
displayed images of compact rank-two modules are closed.

For a single positive prime $p=\ell$, the valuation coordinate is
one and the logarithmic coordinate is zero. If $p\ne\ell$ and
$\ell$ is odd, the logarithm is an isometry on $1+\ell\Z_\ell$, so
\[
 v_\ell(d)=v_\ell(p^{\ell-1}-1)-1.
\]
For $\ell=2$ and $p$ odd, put $\epsilon=(-1)^{(p-1)/2}$.
The logarithm is an isometry on $1+4\Z_2$, and
\[
 v_2(d)=v_2(\epsilon p-1)-2=v_2(p^2-1)-3.
\]
The extra factor two in (D4.2) proves (D4.1). These valuations are
finite because a positive prime is not a root of unity. This proves
openness without discarding the Wieferich cases or the coefficient
prime two. The exact finite-level checks in
\texttt{root\_local\_kummer.py} accompany the general argument.
\end{proof}
